\chapter{Conclusão}\label{cap:conclusao}

Este trabalho investigou a aplicabilidade da coerência \textit{Wavelet} alternativa de Schuck e
Bodmann como descritor dos padrões de escoamento bifásico no arranjo experimental do LETEF,
calculada sobre sinais espaciais de brilho reconstruídos do panorama das imagens da câmera de alta
velocidade, em vez dos sinais dos sensores no domínio do tempo. As seções a seguir sintetizam o que
foi realizado e obtido, discutem o atendimento aos objetivos e apontam os encaminhamentos para
trabalhos futuros.

% =========================================================
\section{Síntese do trabalho}\label{sec:conc-sintese}
% =========================================================

Implementou-se o cálculo da coerência \textit{Wavelet} alternativa a partir das funções de
autocorrelação e correlação cruzada de pares de sinais espaciais, com a transformada \textit{Wavelet}
contínua de Morlet complexa avaliada no atraso nulo, e aplicou-se o método a duas campanhas: a
intra-padrão, sobre os noventa e um pares de ensaios de um mesmo padrão, e a inter-padrão, sobre os
duzentos e oitenta e sete pares de ensaios de padrões distintos. As duas campanhas geraram os campos
de potência, de espectro cruzado e de coerência, além dos perfis de potência média por comprimento
de onda, para cada par de ensaios e nas duas variantes do sinal espacial.

Os resultados mostraram que a coerência \textit{Wavelet} alternativa, tomada como descritor
principal, não descreveu os padrões de escoamento com a precisão esperada: os campos de coerência no
atraso nulo apresentaram-se baixos e irregulares, mesmo entre ensaios de um mesmo padrão. A
caracterização mais consistente veio do perfil de potência média, que concentrou a energia espectral
de cada ensaio em uma faixa de comprimento de onda reprodutível dentro de cada padrão, e foi essa
leitura que melhor evidenciou a repetibilidade intra-padrão. Na comparação inter-padrão, as faixas
dominantes de \textit{churn} e de anular ondulado apresentaram-se parcialmente sobrepostas, de modo
que a separação entre esses padrões mostrou-se sutil.

% =========================================================
\section{Atendimento aos objetivos}\label{sec:conc-objetivos}
% =========================================================

O objetivo geral, investigar a aplicabilidade da coerência \textit{Wavelet} alternativa como
descritor dos padrões de escoamento a partir de sinais espaciais de brilho, foi cumprido: o método
foi implementado, validado com sinais sintéticos e aplicado às duas campanhas sobre os ensaios do
arranjo do LETEF, e sua adequação foi avaliada quanto à repetibilidade intra-padrão e à separação
inter-padrão. A resposta obtida foi que a coerência alternativa, nesse arranjo e com esses sinais,
tem poder discriminante limitado, ao passo que o perfil de potência média se mostrou o descritor
mais promissor.

Os objetivos específicos também foram atendidos: revisou-se a literatura sobre padrões de escoamento
em microcanais e sobre a análise espectral de sinais para sua caracterização; apresentou-se a
fundamentação teórica da transformada \textit{Wavelet} contínua, do espectro de potência, do
espectro cruzado e da coerência nas formulações clássica e alternativa; e descreveu-se a metodologia
de reconstrução do panorama, de pré-processamento e de cálculo da coerência a partir dos sinais
espaciais, aplicada e documentada com os resultados das duas campanhas.

% =========================================================
\section{Trabalhos futuros}\label{sec:conc-futuros}
% =========================================================

O encaminhamento mais imediato será a quantificação estatística das leituras qualitativas obtidas.
Um teste de permutação que tome o ensaio como unidade de replicação poderá avaliar formalmente se a
distância entre assinaturas espectrais separa os padrões além do esperado ao acaso, respeitando a
dependência entre pares que compartilham o mesmo ensaio. Antes da comparação, um branqueamento dos
espectros de potência poderá remover a componente comum de fundo e realçar as diferenças entre
padrões. Será igualmente necessário controlar a influência das condições operacionais, avaliando se a
separação observada persiste quando o fluxo mássico e o fluxo de calor são mantidos comparáveis entre
os padrões. Por fim, a validação do método com sinais sintéticos de fase e razão sinal-ruído
controladas poderá ser ampliada, de modo a delimitar em que condições a coerência alternativa
recupera o acoplamento esperado, esclarecendo os limites observados neste trabalho.
